\documentclass[12pt]{article}
\usepackage[utf8]{inputenc}
\usepackage[T1]{fontenc}
\usepackage{hyperref}
\title{Totela Language Resources for K-12 Curriculum Development}
\author{Compiled by Claude Sonnet 5 (claude-sonnet-5) for Andr\'as Kornai}
\date{}
\begin{document}
\maketitle
\section{Introduction}
Totela (ISO 639-3 ttl; Guthrie K.41) is a poorly described Bantu language of Zambia and possibly Namibia. Wikipedia reports 1,220 speakers in Namibia from the 2010 census. Joshua Project reports the Totela only in Zambia's Western Province, with younger people shifting to Silozi and English. Little material exists.
\section{Primary Educational Resources}
No school materials were found in the search results.
\section{Linguistic and Reference Materials}
\begin{itemize}
\item Wikipedia, Totela language. \url{https://en.wikipedia.org/wiki/Totela_language}
\item Totela numerals, Max Planck Institute numerals database. \url{https://lingweb.eva.mpg.de/channumerals/Totela.htm}
\item A researcher's blog notes the chapter ``Totela K41'' in \emph{The Bantu Languages}, 2nd edition, and a PhD dissertation on Totela. \url{https://blogs.helsinki.fi/crane/?p=235}
\end{itemize}
\section{Religious and Community Resources}
\begin{itemize}
\item Scripture Earth lists a Totela New Testament and audio recordings. \url{https://www.scriptureearth.org/00eng.php?iso=ttl}
\item Joshua Project, Totela in Zambia. \url{https://joshuaproject.net/people_groups/15571/ZA}
\end{itemize}
\section{Digital and Audio Resources}
Scripture Earth (above) links to Global Recordings Network audio and a Bible.com New Testament.
\section{Curriculum Development Context}
The language is thinly documented. A curriculum would need new community-based materials and could draw on the scholarly work above and the scripture texts.
\end{document}